% ---------------------------------------------------------------------------
% Section 4: Generalization and error analysis (Hazel)
% Numbers from: python -m src.q4_analysis --model masked_lr  (results/q4/masked_lr/)
% Paste each answer over the matching \placeholder{...} / table in the template.
% ---------------------------------------------------------------------------

% ===== Q4.1 =====
\question{Q4.1 [5 points].}{Compare training and validation AUROC for your final selected model. Is there a gap? What explanations are plausible?}
On the fixed 8{,}000/2{,}000 split (seed 42) the masked n-gram logistic regression reaches a training AUROC of 0.9998 and a validation AUROC of 0.9949, so the gap is only 0.005. At a 0.5 threshold the error rate rises from 0.6\% on training rows to 2.6\% on validation rows (52 of 2{,}000). The small gap is expected: the model has up to 200{,}000 masked 1--3-gram features for 8{,}000 rows and $C=4$ is weak regularization, so it can nearly memorize the training set, but the structural cues it relies on (blank lines between paragraphs, punctuation and function-word patterns) are shared across the whole training distribution, so they transfer to held-out rows from the same sources. One caution: the validation set is drawn from the same distribution as training, while the test set is shifted towards fiction and poetry (Notebook 01). A leave-cluster-out estimate of 0.988 (Notebook 02) suggests the true test gap is larger than this validation gap.

% ===== Q4.2 =====
\question{Q4.2 [5 points].}{Evaluate performance over at least three text-length bins (define the bins), or another observable subgroup available in your data. Explain why you chose the grouping.}
\begin{tabularx}{\linewidth}{@{}Yrr@{}}\toprule Subgroup / bin & Count & Validation AUROC \\ \midrule
Short: $<150$ words & 404 & 0.9885 \\
Medium: 150--249 words & 1{,}055 & 0.9967 \\
Long: $\geq 250$ words & 541 & 0.9954 \\
\midrule
Format: blank lines between paragraphs & 689 & 1.0000 \\
Format: single line breaks only & 709 & 0.9616 \\
Format: no line breaks & 602 & 0.9838 \\
\bottomrule\end{tabularx}
\par\noindent
We binned by word count because the rubric suggests it and because a bag-of-n-grams model simply has less evidence in a short passage. We added the line-break format because Notebook 02 showed it is the single strongest feature. Short passages are clearly the hardest length bin: 5.9\% errors versus 1.2\% for medium passages. Format matters even more. Passages with blank lines are 99.6\% AI in the validation set and the model gets every one right, while passages with only single line breaks (93\% human) have the lowest AUROC. By topic cluster (results/q4/masked\_lr/q42\_by\_topic.csv), the fiction cluster has the highest error rate, 9.8\%, but only 41 validation rows, which is a concern because about 36\% of the test set falls into that cluster.

% ===== Q4.3 =====
\question{Q4.3 [5 points].}{Analyze at least \textbf{five false positives and five false negatives} from your held-out validation set. Provide example IDs, predicted probabilities, and short failure hypotheses. Do not paste full passages unless essential.}
\begin{tabularx}{\linewidth}{@{}p{0.64in}p{0.70in}p{0.98in}Y@{}}\toprule ID & Error type & P(AI) & Possible reason \\ \midrule
\mbox{1898} & FP & 0.738 & Human psychology article in a formal, textbook register with a single line break before a sub-heading (``Modelling''), a layout the model mostly sees in AI text. \\
\mbox{3139} & FP & 0.689 & Human encyclopedia entry on fluid power: long sentences, lists of technical nouns and no line breaks; with topic words masked, little remains to mark it as human. \\
\mbox{3155} & FP & 0.619 & Human book summary with a double space after a sentence and one paragraph break; neutral, polished prose close to the AI style. \\
\mbox{151} & FP & 0.595 & Short (79 words) movie review with a typo (``berween''); too few n-grams for confident evidence, and the review-style quotes resemble AI film summaries. \\
\mbox{1038} & FP & 0.582 & Human history article with en-dashes and many proper nouns; formal encyclopedic register again. \\
\mbox{5498} & FN & 0.010 & AI recipe written in the human source format: an ingredient list with one item per line and no blank lines, the strongest ``human'' cue in training. \\
\mbox{5010} & FN & 0.027 & AI biography that imitates Wikipedia, including leftover footnote digits (``California1'', ``miniseries.2''), with no line breaks. \\
\mbox{9298} & FN & 0.028 & AI biography of a comedian with quoted titles and a single line break between paragraphs, matching human Wikipedia formatting. \\
\mbox{5929} & FN & 0.072 & AI recipe with messy spacing, missing line breaks between items and a typo (``yollk''), so it looks like a hastily typed human recipe. \\
\mbox{1614} & FN & 0.080 & Short first-person diary-style passage (``I've never felt more alone'') with contractions and no line breaks; informal first-person text is overwhelmingly human in training. \\
\bottomrule\end{tabularx}
\smallnote{Decision threshold 0.5. Examples are the five most confident errors of each type; all 52 validation errors are listed in results/q4/masked\_lr/q43\_all\_errors.csv.}

% ===== Q4.4 =====
\question{Q4.4 [5 points].}{What recurring failure patterns did you identify, and how might you address one of them? Distinguish observed patterns from speculation.}
\textbf{Observed.} (1) Errors are one-sided: 43 of the 52 validation errors are false negatives, AI text predicted human. (2) Almost all of them lack blank lines: 39 of 43 false negatives have either single line breaks or none, and the median false negative has zero newlines per 100 words, against a validation-set mean of 4.1. The model has learned ``blank line between paragraphs means AI'' because that is how the AI training passages are formatted, and it under-reacts when an AI passage copies the layout of its human source (recipes as ingredient lists, Wikipedia-style biographies). (3) The false positives are the mirror image: formal encyclopedic human text with no topic words left after masking. (4) Short passages (under 150 words) have about five times the error rate of medium ones.
\textbf{Speculation.} The test set contains adversarially modified AI passages; if that modification removes paragraph breaks or copies human formatting, the false-negative pattern above will be worse on test than on validation, and the fiction cluster (36\% of test, 9.8\% validation error rate) will amplify it.
\textbf{How to address it.} Reduce the dependence on formatting rather than remove it: train the same model on passages where blank lines are randomly collapsed to single line breaks for half of the training rows (formatting augmentation), or average the masked model with the ``flat'' variant from Notebook 02 that sees no line breaks at all, which was competitive under leave-cluster-out validation and stronger on the fiction cluster. Either option can be tested on the same fixed validation split and is a natural candidate for the Q5 improvement.
